% ==========================================================================
% GENERATED by scripts/split_paper.py. Do not edit this file.
%
% Source: imported_papers/probe-norm/content/body.tex, content/supplement.tex
% Part 11: Literature Survey
%
% Edits belong upstream, in the published paper. To change what the
% chapter includes, change imports/bounded.toml and re-run
% the splitter.
% ==========================================================================
\section{Literature Survey}
\label{app:literature}

\emph{Search.} In October 2026, we searched arXiv and machine-learning and robotics proceedings for work published in 2019--2026 on JEPA world models, latent planning, model-based reinforcement learning, and latent physical or video forecasting. We also searched for predicted-latent readouts and random-dynamics controls, and followed references and citing papers.

\emph{Inclusion.} Studies must predict latents recursively or from actions, apply a readout to those predictions, and report scores or decoded rollouts.

\emph{Exclusion.} We exclude observation-space predictors, probes applied only to encodings, and latents used only inside planners or value functions without a reported readout evaluation. Fifteen retrieved papers fail these criteria, including TD-MPC, TD-MPC2, DreamerV2, SPR, and TCRL.

\emph{Extraction.} We record the readout, its training inputs, loss, recurrent output constraint, trained horizon $H$, evaluated horizons $T$, and reported checks (Table~\ref{tab:literature}). $H$ counts chained prediction steps in the loss, so teacher forcing gives $H=1$ regardless of context length. We resolve missing fields from code where possible and otherwise mark them unstated. Values derived from unstated quantities or inherited architectures are marked inferred. Sub-JEPA and Fast-LeWM retain LeWorldModel's normalized head. VLWM states the same architecture but has no released code. \emph{Normalized input} denotes normalization or re-projection built into the readout. None of these evaluations compares scores with and without that operation.

\emph{Screen.}
%%BEGIN_SURVEY_SCREEN
We screen for quantitative evaluation beyond five times the training horizon with a readout fitted on encodings, no input normalization, and no reported prediction norms. Seven evaluations meet these criteria (Table~\ref{tab:literature}). Five have unconstrained or undocumented recurrent outputs: two Koopman autoencoders use their own decoder, DINO-Foresight uses an auxiliary readout with unconstrained outputs, and Navigation World Models and RDR leave output constraints unstated. The other two have bounded outputs: MILE uses its own decoder, and the imagined-DreamerV3 study uses a separate probe~\citep{schafer2026imagined}. Three LeWorldModel variants report long rollouts only qualitatively and are excluded from this screen, as are evaluations with unknown horizons. Where the model's own decoder reads the predictions, re-projection shows how much of the score the norm causes, and the raw error remains the measure of the model. The screen lists the evaluations that lack this comparison.
%%END_SURVEY_SCREEN

\input{tables/table_literature.tex}
